\scanpage{17}
\mspage{14}

suivantes soient satisfaites

a) Si $I\supset J$, \inserted{\unsure{des segments}}, et si \deleted{$x_0\in\delta I$, \unsure{on}} \inserted{$\unsure{o_I|\sigma_I}$ est un des deux ordres canoniques sur $I$,} \deleted{alors $\delta I\simeq\delta J$ lui fait correspondre l'unique $y_0\in\delta J$ qui soit le plus petit élément de $J$ pour l'ordre} \inserted{alors $o_J:=o|J$ est un des deux ordres canoniques sur $J$, alors $\delta I\simeq\delta J$} associe au plus petit élément de $I$ le plus petit élément de $J$ (et idem pour le plus grand, par conséquent en \deleted{\unsure{en}} remplaçant $o$ par l'ordre opposé)

b) Si $I\in S$, et si $x_0\in\delta I$ est le plus petit élément pour \deleted{l'un des} \inserted{l'ordre} canonique $o$ sur $I$, alors \unsure{le} \deleted{\unsure{à}} $x_0\in\delta I$ correspond la branche en $\unsure{o|x}$ définie par $I$.

Il faut voir que la relation d'équivalence $R$ définie dans
\[
\coprod_{I\in S}\delta I\ \amalg\ \coprod_{x\in L}B_x
\]
par les relations élémentaires \deleted{$\unsure{R_a}\ \unsure{R_b^{\ast}}$} \inserted{\textcircled{a}, \textcircled{b}}
\[
(I,\underset{\substack{\rotatebox{-90}{$\in$}\\\delta I}}{x})
\underset{\substack{\deleted{\unsure{R_a}}\\\text{\textcircled{a}}}}{\sim}
(\unsure{I|J},y)
\quad\Longleftrightarrow\quad
J\subset I\text{ et }x\text{ et }y
\]
se correspondent par la \unsure{prescription} \unsure{a}
\[
(I,\underset{\substack{\rotatebox{-90}{$\in$}\\\delta I}}{x})
\underset{\text{\textcircled{b}}}{\sim}
\left(x,\underset{\substack{\text{branche}\\\text{en }x\text{ définie}\\\text{par }I}}{\underbrace{\beta(I,x)}}\right)
\]
admet exactement deux \uline{classes}, et que

\mspage{16}

\rotatebox{90}{\reflectbox{\unsure{EXTRA STRONG}}}
\rotatebox{90}{\reflectbox{\unsure{EXTRA STRONG}}}
